% ============================================================
%  Curva Característica I-V del Diodo
%  Requiere: tikz, siunitx
%  Librerías TikZ: arrows.meta
% ============================================================

% ── Paleta de colores ──────────────────────────────────────
\definecolor{curvaColor}{RGB}{30, 90, 200}      % azul principal
\definecolor{directaColor}{RGB}{20, 130, 60}    % verde (directa)
\definecolor{inversaColor}{RGB}{180, 70, 20}    % naranja/rojo (inversa)
\definecolor{rupturaColor}{RGB}{160, 20, 60}    % rojo oscuro (ruptura)
\definecolor{fondoDirecta}{RGB}{220, 245, 225}  % verde muy claro
\definecolor{fondoInversa}{RGB}{255, 235, 220}  % naranja muy claro
\definecolor{axisGray}{RGB}{60, 60, 60}

\begin{tikzpicture}[
  >=Latex,
  font=\small,
  thick,
  % Estilo para etiquetas de región
  regionLabel/.style={
    font=\footnotesize\bfseries,
    align=center,
    rounded corners=3pt,
    inner sep=4pt,
    draw=gray!40,
    fill=white,
    opacity=0.95
  }
]

% ─────────────────────────────────────────────────────────────
%  FONDO DE REGIONES  (capa inferior)
% ─────────────────────────────────────────────────────────────
\begin{scope}[on background layer]
  % Región de polarización directa (V > 0)
  \fill[fondoDirecta, opacity=0.6]
    (0.05, -3.4) rectangle (3.1, 4.8);
  % Región de polarización inversa (V < 0)
  \fill[fondoInversa, opacity=0.45]
    (-5.3, -3.4) rectangle (0.05, 4.8);
\end{scope}

% ─────────────────────────────────────────────────────────────
%  EJES
% ─────────────────────────────────────────────────────────────
\draw[axisGray, thick, ->] (-5.4, 0) -- (3.4, 0)
  node[right, font=\small] {$V$ (\unit{V})};
\draw[axisGray, thick, ->] (0, -3.5) -- (0, 5.0)
  node[above, font=\small] {$I$ (\unit{mA})};

% Origen
\node[below left, font=\scriptsize, axisGray] at (0, 0) {$O$};

% Marcas menores en el eje V (escala esquemática)
\foreach \x in {-5,-4,-3,-2,-1,1,2,3} {
  \draw[axisGray, thin] (\x, 0.06) -- (\x, -0.06);
}
\foreach \y in {-3,-2,-1,1,2,3,4} {
  \draw[axisGray, thin] (0.06, \y) -- (-0.06, \y);
}

% ─────────────────────────────────────────────────────────────
%  CURVA CARACTERÍSTICA DEL DIODO
% ─────────────────────────────────────────────────────────────

% 1. Zona de RUPTURA (caída abrupta a -V_BR = -5 V esquemático)
\draw[rupturaColor, very thick, line cap=round]
  (-4.8, -3.2) -- (-4.8, -0.22);

% 2. Corriente inversa de saturación (casi plana, muy pequeña)
\draw[curvaColor, very thick]
  (-4.8, -0.22) -- (-0.35, -0.22);

% 3. Transición suave cerca del origen
\draw[curvaColor, very thick]
  (-0.35, -0.22) ..
  controls (-0.1, -0.22) and (0.02, -0.02) ..
  (0.08, 0);

% 4. Polarización directa sub-umbral (0 → 0.7 V): corriente ≈ 0
\draw[curvaColor, very thick]
  (0.08, 0) -- (1.1, 0.08);

% 5. Polarización directa sobre la rodilla (V > 0.7 V): subida exponencial
%    Recortado al área del diagrama con \clip
\begin{scope}
  \clip (-5.5, -3.5) rectangle (3.3, 4.85);
  \draw[directaColor, very thick, domain=1.1:2.45, samples=80]
    plot (\x, {0.08 + 4.2 * pow(\x - 1.1, 1.75)});
\end{scope}

% ─────────────────────────────────────────────────────────────
%  LÍNEAS DE REFERENCIA PUNTEADAS
% ─────────────────────────────────────────────────────────────

% Voltaje umbral (rodilla) en 0.7 V
\draw[gray!65, dashed, thin] (1.1, 0) -- (1.1, 0.08);

% Voltaje de ruptura
\draw[gray!65, dashed, thin] (-4.8, 0) -- (-4.8, 0.12);

% Nivel de corriente de saturación
\draw[gray!65, dashed, thin] (0, -0.22) -- (-0.35, -0.22);

% ─────────────────────────────────────────────────────────────
%  MARCAS Y ETIQUETAS DE VALORES
% ─────────────────────────────────────────────────────────────

% Tick + etiqueta: 0.7 V
\draw[gray!70, thin] (1.1, 0.07) -- (1.1, -0.07);
\node[below, font=\footnotesize, text=gray!80] at (1.1, -0.1)
  {$\qty{0.7}{V}$};

% Tick + etiqueta: -V_BR
\draw[gray!70, thin] (-4.8, 0.07) -- (-4.8, -0.07);
\node[below, font=\footnotesize, text=rupturaColor] at (-4.8, -0.1)
  {$-V_{BR}$};

% Corriente de saturación inversa
\node[right, font=\footnotesize, text=curvaColor] at (0.08, -0.22)
  {$-I_S$};

% ─────────────────────────────────────────────────────────────
%  ETIQUETAS DE REGIÓN
% ─────────────────────────────────────────────────────────────

% Polarización DIRECTA
\node[regionLabel, text=directaColor] at (2.1, 4.45)
  {Polarización\\Directa};

% Polarización INVERSA
\node[regionLabel, text=inversaColor] at (-2.3, -1.7)
  {Polarización\\Inversa};

% Zona de RUPTURA con flecha
\node[regionLabel, text=rupturaColor] (nRuptura) at (-3.2, -2.8)
  {Zona de\\Ruptura};
\draw[->, rupturaColor!60, thin]
  (nRuptura.north west) -- (-4.6, -1.8);

% ─────────────────────────────────────────────────────────────
%  ANOTACIONES FÍSICAS CLAVE
% ─────────────────────────────────────────────────────────────

% Flecha + nota: barrera de potencial
\draw[<-, gray!60, thin]
  (1.4, 0.2) -- ++(20:1) node[above, align=left, font=\scriptsize, text=gray!70]{Barrera de\\ potencial};

% Flecha + nota: corriente casi nula en inversa
\draw[->, gray!60, thin]
  (-1.8, 0.5) -- (-1.0, -0.2);
\node[above, align=center, font=\scriptsize, text=gray!70]
  at (-1.9, 0.6) {Corriente de\\fuga ($\approx 0$)};

% ─────────────────────────────────────────────────────────────
%  NOTA DE ESCALA (eje V no lineal entre regiones)
% ─────────────────────────────────────────────────────────────
\node[font=\scriptsize, text=gray!60, align=left]
  at (-1, -3.7) {\textit{Nota: escala del eje $V$ no lineal (esquemática)}};

\end{tikzpicture}
